\usepackage{times}
\usepackage[T1]{fontenc}
\usepackage{amsmath,amssymb,mathtools}
\usepackage{xcolor}
\usepackage{algorithm}
\usepackage[indLines=true,noEnd=true]{algpseudocodex}
\usepackage{booktabs,tabularx}
\definecolor{commentink}{gray}{0.48}
\definecolor{slotink}{RGB}{35,66,86}
\algrenewcommand\algorithmiccomment[1]{\textcolor{commentink}{\#\ #1}}
\algrenewcommand\alglinenumber[1]{\scriptsize\textcolor{commentink}{#1}}
\algrenewcommand\algorithmicindent{1em}
\newcommand{\slot}[1]{\textcolor{slotink}{\text{\sffamily\bfseries #1}}}
\newcommand{\fn}[1]{\operatorname{\textnormal{\textsc{#1}}}}
\newcommand{\Unif}{\operatorname{Unif}}
\newcommand{\argmin}{\operatorname*{arg\,min}}
\newcommand{\argmax}{\operatorname*{arg\,max}}
\newcommand{\clip}{\operatorname{clip}}
\newcommand{\ind}{\mathbf{1}}
\newcommand{\io}[2]{\noindent\textbf{#1}\enspace #2\par\smallskip}
\newcommand{\algnotes}[1]{\par\smallskip\noindent{\small #1}\par}
\newcommand{\slotsnote}{\textcolor{slotink}{\textsf{\textbf{Bold sans-serif}}} marks an evolvable operator slot.}
\floatname{algorithm}{Algorithm}
